\centerline{\bf \Huge ГМТ}
\centerline{\bf \Huge }

\begin{flushright}
        \scriptsize{Я только узнал, что ГМТ на английском это просто locus}
\end{flushright}

\textit{Геометрическое место точек} (сокращённо ГМТ), обладающих некоторым свойством, --- это фигура, состоящая из всех точек, для которых выполнено это свойство.

Решение задачи на поиск ГМТ должно содержать доказательство того, что

a) точки, обладающие требуемым свойством, принадлежат фигуре $\Phi$, являющейся ответом задачи;

б) все точки фигуры $\Phi$ обладают требуемым свойством.


Три важнейших ГМТ:

a) ГМТ, равноудалённых от точек $A$ и $B$, является серединным перпендикуляром к отрезку $A B$ \textbf{(Доказать!)};

б) ГМТ, удалённых на расстояние $R$ от данной точки $O$, является окружностью радиуса $R$ с центром $O$ (Это определение окружности, доказывать не надо);

в) ГМТ, из которых данный отрезок $A B$ виден под данным углом, является объединением двух дуг окружностей, симметричных относительно прямой $A B$ \textbf{(Доказать!)}.

\problem{1}
а) Докажите, что ГМТ, равноудалённых от сторон угла меньше развёрнутого, является его биссектриса.

\noindent
б) Найдите все точки равноудалённые от двух данных пересекающихся прямых.

\noindent
в) Докажите, что три биссектрисы треугольника пересекаются в одной точке.

\problem{2}
Точка $O$ --- середина отрезка $M K$. Известно, что $A M<B M, A K<B K$. Докажите, что $A O<B O$.

\problem{3}
Дан выпуклый шестиугольник, никакие стороны которого не параллельны. Сколько существует точек, равноудалённых от трёх его несмежных сторон?

\problem{4}
а) Дан равносторонний треугольник $A B C$. Найдите множество таких точек $M$, что $B M<AM, CM$.

\noindent
б) Дан квадрат $A B C D$. Где находятся все такие точки $M$, что $A M \geqslant B M \geqslant C M \geqslant D M$ ?


\newpage

\hspace{3mm}

\problem{5}
Один из углов треугольника равен $120^{\circ}$. Докажите, что треугольник, образованный основаниями биссектрис данного, --- прямоугольный. (Указание. Рассмотрите внешний угол)

\problem{6}
В четырёхугольнике $A B C D$ внешний угол при вершине $A$ равен углу $\angle B C D, A D=$ $C D$. Докажите, что $B D$ - биссектриса угла $A B C$.